% A1Completeness20261024: resistor-integral
下面把最后的积分算清，避免把电路答案藏在“标准积分”中。对$a>|b|$，设
\[
 J(a,b)=\frac1{2\pi}\int_{-\pi}^{\pi}\frac{\dd k}{a-b\cos k}.
\]
令$t=\tan(k/2)$，则$\cos k=(1-t^2)/(1+t^2)$、$\dd k=2\dd t/(1+t^2)$。
$k$从$-\pi$走到$\pi$时，$t$从$-\infty$走到$+\infty$，故
\[
 J=\frac1\pi\int_{-\infty}^{\infty}
 \frac{\dd t}{(a-b)+(a+b)t^2}.
\]
因$a-b,a+b$都为正，再令$u=t\sqrt{(a+b)/(a-b)}$：
\[
 J=\frac1{\pi\sqrt{a^2-b^2}}
 \int_{-\infty}^{\infty}\frac{\dd u}{1+u^2}
 =\frac{\arctan u|_{-\infty}^{\infty}}{\pi\sqrt{a^2-b^2}}
 =\frac1{\sqrt{a^2-b^2}}.
\]
\begin{qualityfigure}
\begin{tikzpicture}[>=Latex,font=\small]
 \begin{scope}
  \draw[->](-2.2,1)--(2.3,1)node[right]{$k$};
  \foreach\x/\lab in{-1.8/-\pi,0/0,1.8/\pi}{\draw(\x,.92)--(\x,1.08);\node[above]at(\x,1.1){$\lab$};}
  \draw[->](-2.2,-.7)--(2.3,-.7)node[right]{$t$};
  \node[below]at(-1.8,-.8){$-\infty$};\node[below]at(0,-.8){$0$};\node[below]at(1.8,-.8){$+\infty$};
  \foreach\x in{-1.8,0,1.8}{\draw[->,DeepBlue](\x,.65)--(\x,-.35);}
  \node at(0,-1.7){$t=\tan(k/2)$};
 \end{scope}
 \begin{scope}[shift={(6.2,-.4)}]
  \draw[->](-2.1,0)--(2.3,0)node[right]{$u$};
  \draw[->](0,0)--(0,1.65)node[above]{$(1+u^2)^{-1}$};
  \fill[PaleBlue,domain=-2:2,samples=80](-2,0)--plot(\x,{1/(1+\x*\x)})--(2,0)--cycle;
  \draw[DeepBlue,thick,domain=-2:2,samples=80]plot(\x,{1/(1+\x*\x)});
  \node at(0,-.65){全实轴曲线下的面积为$\pi$};
 \end{scope}
\end{tikzpicture}
\figcap{半角代换将一个周期的积分化为全实轴上的反正切积分；右图的两端尾部延伸到无穷远。}\label{advA1:fig:completeness-resistor-integral}
\end{qualityfigure}
图~\ref{advA1:fig:completeness-resistor-integral}左侧的三个对应点固定了积分区间，右侧面积由反正切的两个端点极限得到。本题取$a=r^{-1}+2R^{-1}$、$b=2R^{-1}$，电阻为正保证$a-b=r^{-1}>0$，代入得
